\DocumentMetadata{
  lang          = en-US,
  pdfversion    = 2.0,
  pdfstandard   = ua-2,
  tagging       = on,
  tagging-setup = {math/setup=mathml-AF, math/alt/use}
}
\documentclass[12pt]{article}
\usepackage[T1]{fontenc}
\usepackage[utf8]{inputenc}  
\usepackage[american]{babel}
\usepackage[authordate,backend=biber]{biblatex-chicago}
\addbibresource{references.bib}
\usepackage{newpxtext}       
\usepackage{csquotes}%
\usepackage{newpxmath}       
\usepackage{graphicx}       
\usepackage{geometry}
\usepackage{float}
\usepackage{lipsum}
\usepackage{mathtools}
\usepackage{setspace}
\usepackage{parskip}
\usepackage{hanging}
\usepackage{indentfirst}
\usepackage{authblk}
\usepackage{threeparttable}
\usepackage{booktabs}
\usepackage[dvipsnames]{xcolor}
\definecolor{BingGreen}{RGB}{0,93,64}
\definecolor{AccentGrey}{HTML}{5A5C5B}
\definecolor{AccentGreen}{HTML}{417c2a}
\usepackage{hyperref}
\hypersetup{
  pdftitle           = {PLSC 389U: Internal Conflict},
  pdfauthor          = {Alexander D. Perdue},
  pdfkeywords        = {Social movements, collective contention, protest, vigilantism, political violence},
  colorlinks         = true,
  linkcolor          = BingGreen,
  urlcolor           = BingGreen,
  citecolor          = AccentGreen,
  pdfdisplaydoctitle = true,
}
\DeclareCiteCommand{\bibitemcite}
   {\setcounter{maxnames}{99}\setcounter{minnames}{99}}
  {\usedriver{\frenchspacing}{\thefield{entrytype}}}
  {\par}
  {}

\title{\textbf{PLSC 389U: The Politics of Internal Conflict}}
\author{Alexander (Alex) D. Perdue}
\date{Spring 2027}

\begin{document}

\maketitle

\section*{Course Description}
Not every incident of collective contention—organized contestation over politics between groups and/or individuals—ends in, or occurs during, a civil war. Organized contention might be mundane: general social movements for change, changes in law enforcement, or campaigns and elections. We will look beyond these everyday forms of politics to examine collective contention in its more violent forms short of civil war, focusing on modes of politics such as vigilantism, mass violence, and finally, anti-state movements. This course gives students an introduction to the study of internal contention and contentious politics, exploring how social scientists think about the range of violent outcomes that occur during ostensible peace.

\section*{Course Information}
\begin{itemize}
    \item Class Time: Mondays and Wednesdays 5:00-6:30
    \item Class Location: CW 112
    \item My Office: LNG 279
    \item My Office Hours: Mondays and Wednesdays 3:30-4:45
    \item My Email: \textbf{\textcolor{BingGreen}{\href{aperdue2@binghamton.edu}{aperdue2@binghamton.edu}}} 
\end{itemize}

\subsection*{Course Logistics}
\par{Once the course opens, please feel free to email me at any time at the email provided above. 
\begin{itemize}
    \item Emails sent on weekdays should receive a reply within 24 hours.
    \item Emails sent on the weekend should receive a reply within 48 hours.
    \item Requests to meet should be given at least 24 hours before you want to meet unless you are coming to office hours, which require no prior notification.
\end{itemize}
Please include [PLSC 389U] in your subject line when emailing me just to ensure your email doesn't get shuffled down through my inbox or ignored.}

\section*{Learning Outcomes}
This upper-level political science course fits into several requirements. The first is the General Education \textit{\textbf{N}} requirement. This class contributes toward that requirement by encouraging students to:
\begin{itemize}
    \item Demonstrate knowledge of the major methods, approaches, and issues (and their interrelationships) in political science.
    \item Develop a basic understanding of methods used by social scientists to understand social phenomena including, but not limited to: theory and hypothesis development, measurement, case selection and data collection, and the evaluation of evidence in one or more disciplines in the social sciences.
    \begin{itemize}
        \item In this course, we will pay special attention to issues of measurement, case selection, and the evaluation of evidence.
    \end{itemize}
\end{itemize}

For those of you who are Political Science majors, this course also fits into several other requirements regardless of the specialization or ``track'' you're following, making progress toward your degree requirements as:
\begin{itemize}
    \item \textit{1} of the required \textit{10} PLSC Courses.
    \item \textit{1} of the required \textit{6} 300 or 400-level courses in the major field.
    \item \textit{1} course in Comparative Politics which counts toward the \textit{3}-course requirement in the Comparative Politics subfield for all majors.
    \item \textit{1} course in Comparative Politics which counts toward the \textit{4} course requirement in the Comparative Politics subfield for majors on the Global and International Affairs Track.
\end{itemize}

Additionally, while these are not exhaustive, and I certainly hope you can take more from this course than just these, these are some outcomes I hope we can achieve both individually and as a group for this course:
\begin{itemize}
    \item Study, identify, and reflect on the importance of contentious politics for the study of politics.
    \item Identify and understand processes important for the development of, maintenance of, and endings of contentious collectives and moments of contention.
    \item Describe, evaluate, and explain one or more episodes of contention currently in operation.
    \item Develop an understanding of how processes outside of the traditional institutions framework, at the individual level, might impact the political calculus of everyday citizens.
\end{itemize}

\section*{Academic Honesty}
\subsection*{AI}
\par{\textbf{The use of generative AI tools or apps for assignments in this course, including tools like ChatGPT, Claude, Gemini, and other AI writing or coding assistants, is prohibited.}}

\par{Simply, when you write, I want you to think boldly, and I want to hear \textit{your} thoughts on collective contention, broadly this course's central concern, not someone else's. While outsourcing your reading and reflecting to machines can be a way to save the time that you would spend engaging with material, it is a poor substitute to actually doing the work. Yes, that work may not always be easy (particularly when thinking about the work of Charles Tilly!), it may also not be the most fun for those of you who aren't intending a career in academic political science. However, thinking through your problems and writing your own material is not just best practice, it grants you the ability to grow as a writer, a thinker, and a student, and it gives me the ability to respond to \textit{you}. As your professor, I want to be able to work with you to discover your interests, sharpen your skills, and think through big problems---using a heuristic like AI doesn't help me, or you, do that!}

\par{\textbf{The penalty for confirmed AI use is a 0 on the original assignment and an opportunity to resubmit.} If the student elects to resubmit, the grades of the two assignments will be averaged.}

\subsection*{Plagiarism}
\par{You will not be helped on your assignments by using others' ideas and representing them as your own. This course is, at least in part, about \textit{your} ability to draw connections and make claims. Part of the reason we draw connections and make claims is to make us identifiably us and generate new and exciting ideas! However, when it comes to others' ideas, it's perfectly fine if we attribute others' ideas to their progenitors, not so much otherwise. Thus, if TurnItIn's tools lend me reasonable suspicion to suspect a violation or one or more points governing \href{https://catalog.binghamton.edu/content.php?catoid=2&navoid=66#student-academic-honesty-code}{\textbf{plagiarism, fabrication, or cheating}} in the Binghamton University Student Handbook, I will pursue a violation of the Student Academic Honesty Code with attendant procedures therein.}

\par{\textbf{The penalty for plagiarism is a 0 on the original assignment and an opportunity to resubmit.} If the student elects to resubmit, the grades of the two assignments will be averaged.}

\section*{Accommodations}
\par{If you are requesting academic accommodations for any reason, please do so by registering with the Services for Students with Disabilities Office as early as possible. They are the avenue and adjudicator for disability-related accommodations on campus and they should be your first line of contact to obtain, adjust, or dispute the implementation of your academic accommodations. Their office is in University Union 119 and can be contacted either by phone at (607) 777-2686 or by email at \href{ssd@binghamton.edu}{\textbf{ssd@binghamton.edu}}.}

\par{After I receive your letter, I will reach out to you individually to meet, go over your accommodations, and determine how best to support your needs in the classroom this upcoming semester!} 

\section*{Resources}
\par{Below are a list of potentially useful resources both for this course and more generally during your academic, social, and physical lives at Binghamton. This is not an exhaustive resource list, but it does contain some common ones that might prove to be useful in this setting.}

\subsubsection*{Institutional Resources}
\begin{itemize}
\item Dean of Students Office: 607-777-2804, University Union West 205
\item Office of International Student \& Scholar Services: 607-777-2510, Old Champlain 142
\item Office of Veteran and Military Services: 607-777-3668, University Union West 204 
\item Harpur Advising: 607-777-6305, Old Champlain 110
\item Writing Center: Library North 2411
\end{itemize}

\subsubsection*{Mental, Physical, and Emotional Health Resources}
\begin{itemize}
\item Decker Student Health Services Center: 607-777-2221, behind Johnson Residence Hall in the Dickinson Community
\item University Counseling Center: 607-777-2772, Old O'Connor 264
\item Violence, Abuse, and Rape Crisis Center: Old Johnson 3rd Floor, \textit{Crisis Victim} Assistance Line: 607-777-4256
\item CARE Team: 607-777-2804, University Union West 208G
\end{itemize}

\subsubsection*{Title IX Responsibility---Mandatory Reporter Status}
\par{In addition to the services above, I need to highlight one very important issue you should be aware of before you disclose anything to me or any other office on campus. Some resources, including me, are \href{https://www.binghamton.edu/services/title-ix/about/index.html#:~:text=A%20PRIVATE%20person%2Foffice%20is,to%20the%20greatest%20extent%20possible}{\textbf{Private}} resources, while others are \textbf{\color{AccentGreen}{Confidential}} resources. As a \textbf{\color{BingGreen}{Private}} resource, I am obligated to report any information I receive which concerns sexual assault, harassment, or violence, dating or domestic violence, and/or stalking to the Title IX team on campus.}

\par{If you report something in one of those categories to me or something which could plausibly involve one or more of those, I will report it to the Title IX office. From there, they will likely contact you by email to provide you with on- and off-campus resources and the option to speak with someone about your rights and available options.}

\section*{Assessment}
\subsection*{Grade Scale}
Unless announced to the class otherwise, the following scale for grades will apply in this course:

\tagpdfsetup{table/header-rows={1}}
\begin{center}
\begin{table}[H]
\centering
\caption{Grade scale showing letter grades and their corresponding percentage ranges.}
\begin{tabular}{ccc}
    \toprule
     & Lower Bound & Upper Bound \\
    \midrule
    \textbf{A}  & 94   & 100   \\
    \textbf{A-} & 89.5 & 93.99 \\
    \textbf{B+} & 86.5 & 89.49 \\
    \textbf{B}  & 83.5 & 86.49 \\
    \textbf{B-} & 79.5 & 83.49 \\
    \textbf{C+} & 76.5 & 79.49 \\
    \textbf{C}  & 73.5 & 76.49 \\
    \textbf{C-} & 69.5 & 73.49 \\
    \textbf{D}  & 59.5 & 69.49 \\
    \textbf{F}  & 0    & 59.49 \\
    \bottomrule
\end{tabular}
\end{table}
\end{center}
\vspace{1em}
\par{It is my intention to turn things around as soon as I possibly can relative to the due date, respecting that I am a PhD candidate and have professional obligations. That being said, I will post grades as soon as they are available on Brightspace with feedback (if applicable).}

\subsection*{Deliverables}
\par{Please note that the references for the pre-prepared written assignments you turn in for this class need to conform with the \textbf{\href{https://connect.apsanet.org/stylemanual/}{American Political Science Association style manual}}. While I know many of you will have facility with other style manuals, APSA style is the political science standard. As a result, each pre-pared assignment you turn in must include proper references in APSA formatting.

To aid students, I have marked those assignments with three bold asterisks \textbf{***} before the percentage allocation of the grade to which this applies.}

\par{In this course, you will be graded through the following 4 grade categories.}

\subsubsection*{Quizzes --- 30\%, 5\% Each}
\par{Approximately every other week, students will take an in-class quiz which will cover the material presented since the previous quiz and the readings for each week. While the format of quizzes is at my sole discretion, quizzes might contain a mixture of true/false, multiple choice, and short answer type questions.}

\subsubsection*{Midterm Exam --- 25\%}
\textbf{March 17 during Class Time}
\par{Students will complete a midterm exam right after the midway point of the semester. The format of the exam is at my sole discretion but may include any number of true/false, multiple choice, and short answer questions. Specifics about the format of the exam will be discussed closer to time.}

\subsubsection*{Episodes of Global Contention Presentation \textbf{***} --- 20\%}
\textbf{Presentations Due: Sunday, May 2, 2027 at 11:59 PM EDT}
\par{As the penultimate assignment for the course, you and (a) randomly assigned partner(s), will investigate, write about, and present on some non-civil war episode of political contention that has occurred in the last ten years (e.g. vigilantism in South Africa, anti-gang violence in Brazil) from a list of possible choices I will circulate closer to selection time\footnote{Students \textit{will} have the opportunity to choose something not on the list but to do so they must come to me and defend the decision. Why is it so important to pick this case over other options?}. The goal of this work is to encourage students to think about how developments, technologically (AI, new social media, etc.), politically (the Ukraine War, the Iran War, COVID, etc.), and socially (increasing age polarization, the rising salience of youth-led political mobilization) might influence collective contention and internal conflict. You will need to discuss the following in your presentation:}
    \begin{itemize}
    \item Emergence (20\%) --- Why did your conflict begin where it did? How did it spread between the units it did or why is it so localized?
    \item Operations (20\%) --- What has happened within your conflict? Did it become linked with other conflicts or with social movements? 
    \item Conformance With Theory (25\%) --- Did the episode of contention fit the models prescribed by the theories we discussed in class? If not, how did they differ? If so, how?
    \item Consequences (25\%) --- What was the end result of the episode of contention? Do any of the previous points point to what happened and what the consequences were for participants and the societies they lived in?
    \item Style/Spelling/Grammar (10\%) --- Write well, cite well. Make sure your references are in the appropriate style.
    \end{itemize}

\subsubsection*{Final Exam --- 25\%}
\textbf{Final Exam Time as Determined by the Registrar}
\par{Students will also complete a final exam at the end of the semester. The format of the exam is at my sole discretion but may include any combination of true/false, multiple choice, and short answer questions. Students will also complete at least one longer written response as part of the final exam. Specifics about the format of the exam will be discussed closer to time.}

\par{Of note on this, I will communicate to students when our final exam time is and where it will occur when those facts are made available to me. I will continue to remind students of it as we get closer to time.}

\subsection*{Late Work Policy}
\par{The late work policy in this course will differ based on assignment:
\begin{itemize}
    \item For items which are denoted as needing to be turned in at a specific time already (i.e. the Episodes of Global Contention project), a 20\% per day late penalty will apply.
    \item For items which occur during our regularly scheduled class meeting and will be administered rather than participated in (i.e. the quizzes or the midterm exam), students should contact me if they miss or plan to miss the appointed hour. You will make up the assignment at a time mutually agreeable to us both. 
    \item Making up, doing late, or doing early the final exam will only occur according to the Student Handbook. Once exam schedules are released, it is the student's responsibility to contact me about the need for alternate arrangements, which will only be given according to the Student Handbook (serious personal illness, urgent and serious family matter such as death or severe illness, 3 or more exams scheduled within a 24-hour period where this course is the smallest enrollment course) and only after I receive documentation of the circumstance requiring the work-around.
\end{itemize}}

\section*{A Note on the Material}
\par{As part of this class, you will read material that, at times, might be upsetting. While I have selected materials that try to minimize the potentially upsetting parts of contentious politics, I cannot guarantee that all items might be clear of any upsetting material given the subject matter drawn from. However, I am not the ultimate broker of what might be a trigger, upsetting, or potentially harmful for a student. \textbf{Therefore, I strongly encourage students to consider whether they are capable of engaging with the material and if necessary, please do reach out to me if you experience or anticipate difficulties or a qualified mental health professional if you need additional support.}}

\section*{Required Texts}

\par{Students should acquire the following texts for this course. I would \textit{prefer} students obtain physical copies but acknowledge that texts are expensive and so I am format agnostic. All three can be read online through the libraries if students choose.

\bibitemcite{Cunningham2013}

\bibitemcite{Smith2019}

\bibitemcite{Tilly2003}

Most other items for this course are either journal articles or working papers and thus will be available via the libraries. Therefore, you may obtain the materials by:}
\begin{enumerate}
    \item Copying and pasting the title of the article into the libraries' search engine. It will pop up. Click on a/the link to the article. If prompted, log into your Binghamton student account. You will be able to access a full copy of the text and ought to be able to download it.
    \item Clicking on the doi link. It will open onto the landing page for the article (the same one you would eventually encounter with more steps if you elect the first option). It might prompt you to log in to your student account when you try to download. Do so and you should have the same level of access you otherwise ordinarily would. If not, follow the steps in step (1).
\end{enumerate}

Where items are not available in one of those formats, that is for the one item not yet published, it is available on Brightspace, something I have noted [in brackets].

\section*{Course Schedule}

\textbf{Note}: ``Smith (2019), Introduction and 1'' does NOT mean ``read the introduction of Smith (2019) and page 1''. It means ``read the introduction of Smith (2019) and \textit{chapter} 1.''

\subsection*{Week 1---January 19-22, 2027: Course Introduction}

\textbf{Required Readings}:
\begin{itemize}
    \item Wednesday: Tilly (2003), 1
\end{itemize}

\subsection*{Week 2---January 25-29, 2027: Why Might People Be Violent?}

\textbf{Required Readings}:
\begin{itemize}
    \item Monday: Tilly (2003), 2
    \item Wednesday: Tilly (2003), 10

\end{itemize}

\subsection*{Week 3---February 1-5, 2027: How Might They Respond to the Impetus?}
\textbf{Quiz 1 Will Occur On Wednesday}

\textbf{Required Readings}:
\begin{itemize}
    \item Monday: Tilly (2003), 3
    \item Wednesday: Smith (2019), 7
\end{itemize}

\textbf{Administrative Deadlines}:
\begin{itemize}
    \item Add/Drop (No Penalty) Deadline --- Monday, February 1, 2027 at 11:59 PM EST
\end{itemize}

\subsection*{Week 4---February 8-12, 2027: Vigilantism I---Defining A Practice}

\textbf{Required Readings}:
\begin{itemize}
    \item Monday: Smith (2019), Introduction and 1
    \item Wednesday: \bibitemcite{Nussio2024}
\end{itemize}

\subsection*{Week 5---February 15-19, 2027: Vigilantism II---Historicizing A Practice}
\textbf{Quiz 2 Will Occur This Week}

\textbf{Required Readings}:
\begin{itemize}
    \item Monday: Smith (2019), 2
    \item Wednesday: \bibitemcite{Perdue2027} [On Brightspace]
\end{itemize}

\subsection*{Week 6---February 22-26, 2027: Vigilantism III---Why Keep Vigil?}

\textbf{Required Readings}:
\begin{itemize}
    \item Monday: Smith (2019), 4 and 5
    \item Wednesday: \bibitemcite{Obert2018}
\end{itemize}

\newpage{}
\subsection*{Week 7---March 1-5, 2027: Vigilantism IV---How Do You Bring Down A Vigilante?}
\textbf{Pairings for Episodes of Global Contention Presentation}\\
\textbf{Quiz 3 Will Occur This Week}

\textbf{Required Readings}:
\begin{itemize}
    \item Monday: Trelease ([2022] 1971), 24 and 25
    \item Wednesday: Smith (2019), 8 and 9
\end{itemize}

The chapters you need to read for Monday are from the following book, which is available online via the libraries:\\

\begin{hangparas}{.25in}{1}
\bibitemcite{Trelease1971}
\end{hangparas}
\vspace{0.25em}

\subsection*{Spring Break---March 8-12, 2027}

\subsection*{Week 8---March 15-19, 2027: Midterm Exam}

\textbf{Weekly Schedule}:
\begin{itemize}
    \item Monday: Exam Review Session
    \item Wednesday: Midterm Exam
\end{itemize}

\subsection*{Week 9---March 22-26, 2027: Interpersonal Violence I---Attacks}

\textbf{Required Readings}:
\begin{itemize}
    \item Monday: Tilly (2003), 8
    \item Wednesday: \bibitemcite{Millington2018}
\end{itemize}

\subsection*{Week 10---March 29-April 2, 2027: Interpersonal Violence II---Brawls}
\textbf{Quiz 4 Will Occur This Week}

\textbf{Required Readings}:
\begin{itemize}
    \item Monday: Tilly (2003), 7
    \item Wednesday: \bibitemcite{Batto2020}
\end{itemize}

\textbf{Administrative Deadlines}:
\begin{itemize}
    \item Grade Option Change (P/F) Deadline --- Monday, March 29, 2027 at 4:30 PM EDT
    \item Add/Drop (W) Deadline --- Monday, March 29, 2027 at 11:59 PM EDT
\end{itemize}

\subsection*{Week 11---April 5-9, 2027: Interpersonal Violence III: The Riot}

\textbf{Required Readings}:
\begin{itemize}
    \item Monday: Tilly (2003), 4 and 5
    \item Wednesday: \bibitemcite{Poitras2020}
\end{itemize}

\newpage{}
\subsection*{Week 12---April 12-16, 2027: Contentious Collectives I: The Road to Mainstreaming}
\textbf{Quiz 5 Will Occur This Week}

\textbf{Required Readings}:
\begin{itemize}
    \item Monday: \bibitemcite{Aidt2019}
    \item Wednesday: Cunningham (2013), Introduction and 1
\end{itemize}

\subsection*{Week 13---April 19-23, 2027: Contentious Collectives II: Movement Openings}
\textbf{No Class Wednesday due to Passover Break}

\textbf{Required Readings}:
\begin{itemize}
    \item Monday: 
    \begin{itemize}
        \item Cunningham (2013), 2
        \item \bibitemcite{McVeigh2001}
    \end{itemize}
    \item Wednesday: NO CLASS
\end{itemize}

\newpage{}
\subsection*{Week 14---April 26-30, 2027: Contentious Collectives III: The Rise and Fall of Movements}
\textbf{Quiz 6 Will Occur This Week}

\textbf{Required Readings}:
\begin{itemize}
    \item Monday: Cunningham (2013), 4 and 5
    \item Wednesday: Course Wrap-up:  Cunningham (2013), 6 and 7
\end{itemize}

\textbf{Deliverables Due}:
\begin{itemize}
    \item Sunday, May 2, 2027: Episodes of Global Contention Project Presentation, 11:59 PM EDT
\end{itemize}

\subsection*{Week 15---May 3-7, 2027: Final Presentations}
\begin{itemize}
    \item Monday: Presentations
    \item Wednesday: Presentations
\end{itemize}
\end{document}
